\begin{tabular}{lrrrrrr}
\toprule
Variable & Mean & SD & p10 & Median & p90 & $N$ \\
\midrule
\multicolumn{7}{l}{\emph{Wartime exposure}} \\
\quad Army/AAF deaths & 97 & 315 & 11 & 43 & 157 & 3,073 \\
\quad Fatal burden $B_i$ (per 1{,}000 residents) & 2.74 & 9.26 & 1.35 & 2.30 & 3.63 & 3,073 \\
\quad War supply contracts (\$000s) & 57,980 & 406,004 & 0 & 0 & 43,911 & 3,073 \\
\quad Publicly financed war plant (\$000s) & 4,882 & 26,995 & 0 & 0 & 3,891 & 3,073 \\
\quad Military installations (\$000s) & 3,197 & 12,214 & 0 & 0 & 7,903 & 3,073 \\
\multicolumn{7}{l}{\emph{Prewar characteristics, 1940}} \\
\quad Population & 42,847 & 145,241 & 5,822 & 18,698 & 69,078 & 3,073 \\
\quad Urban share & 0.229 & 0.247 & 0.000 & 0.181 & 0.597 & 3,073 \\
\quad Black population share & 0.106 & 0.178 & 0.000 & 0.012 & 0.394 & 3,073 \\
\quad Manufacturing employment share & 0.028 & 0.038 & 0.000 & 0.012 & 0.080 & 3,073 \\
\quad Log population growth 1930--40 & 0.051 & 0.148 & -0.106 & 0.045 & 0.199 & 3,073 \\
\multicolumn{7}{l}{\emph{Postwar outcomes}} \\
\quad Population 1950 & 49,043 & 171,132 & 5,556 & 18,734 & 80,602 & 3,072 \\
\quad Population 1970 & 65,548 & 232,065 & 5,014 & 18,780 & 111,825 & 3,065 \\
\quad Manufacturing share 1960 & 18.85 & 12.90 & 3.20 & 17.10 & 37.80 & 3,071 \\
\quad Manufacturing share 1970 & 21.85 & 13.57 & 4.10 & 21.00 & 41.00 & 3,066 \\
\quad Median family income 1950 (\$) & 2,240 & 865 & 1,108 & 2,311 & 3,319 & 3,071 \\
\quad Median family income 1970 (\$) & 7,460 & 1,853 & 5,158 & 7,398 & 9,816 & 3,066 \\
\quad Homeownership rate 1970 (\%) & 70.9 & 8.1 & 60.6 & 72.0 & 79.9 & 3,066 \\
\bottomrule
\end{tabular}